\Needspace{15\baselineskip}
\section{The Preparation of the Gifts}

The Liturgy of the Eucharist begins with preparation, not consecration. Bread and wine are brought to the altar; the altar and vessels are made ready; the Church blesses the Creator for gifts that will become the sacrament of Christ's sacrifice; alms embody the communion that the Eucharist demands; priest and people place themselves beneath a proper Prayer over the Offerings. These actions correspond to Christ's taking bread and the cup before giving thanks, breaking, and giving them (GIRM 72). They have real offertory meaning, but they neither repeat Calvary nor anticipate as completed what the Eucharistic Prayer will effect.

The movement is sometimes obscured in opposite ways. It can be treated as an intermission filled by music, as though the Word had ended and the Eucharist had not begun. It can also be overloaded as though every object, intention, and human effort were already sacramentally offered in the same manner as Christ's Body and Blood. The rite supplies a more exact grammar. Created matter and human work are received; ecclesial charity is enacted; the ordained Priest performs his proper ministry; the baptismal people consciously join the offering; and all pass into the great prayer in which the Church, through Christ and in the Holy Spirit, offers the unbloody memorial of the one sacrifice.

\ordermovement{Preparation of the Altar, Procession, Gifts, Collection, and Chant}{Creation and Charity Moving toward the Altar}{Fixed altar preparation; presentation of bread and wine praiseworthy; gifts for the poor or the Church admitted but kept from the Eucharistic table; chant proper or lawfully selected}{GIRM 44, 72--74, 139--140}

After the Prayer of the Faithful, ministers prepare the altar. The corporal, purificator, chalice, pall if used, and Missal are arranged there as required; other vessels may be brought according to need. Bread and wine for the Eucharist are then presented. It is praiseworthy for members of the faithful to bring them, and the Priest or Deacon receives them at a suitable place before they are carried to the altar (GIRM 73, 139--140). Money or other gifts for the poor or the Church may accompany this presentation, but they are placed somewhere suitable away from the Eucharistic table. That separation is theological as well as practical: alms are true gifts of ecclesial charity, while only the appointed bread and wine are the matter to be consecrated.

The action gathers several relations without collapsing them. The altar is the center to which the created gifts move. The corporal marks the place at which the sacramental species will rest. Vessels serve the holy action without themselves becoming its object. Persons carry the gifts, other persons receive and arrange them, and the assembly may sing while the collection proceeds. No one item represents ``the people'' to the exclusion of the rest. The human procession, the products of earth and cultivation, the shared material support of worship and the poor, and the ordered ministries together disclose a Church that brings to God only what it first received.

Scripture gives this receptive offering its grammar. David's prayer over the materials for the Temple refuses possessiveness: wealth comes from God, and the people give from what God's hand has already supplied (1~Chr.~29:10--18). Melchizedek's bread and wine can be read within the Church's larger typology of Christ's priesthood (Gen.~14:18; Heb.~7), but the Genesis scene is not a rubric for the modern procession. Paul's appeal that Christians present their bodies as a living sacrifice extends worship into obedient life (Rom.~12:1--2); it does not identify every human aspiration with Eucharistic matter. Hebrews keeps all Christian self-offering beneath Christ's once-for-all self-gift (Heb.~10:1--18). The gifts move toward an altar whose sacrificial truth is his, not theirs.

The action is securely ancient even though its current ceremonial is not a reconstruction. Justin Martyr describes bread and a cup of wine mixed with water being brought to the presider, and he places voluntary giving for the needy in the same account of Sunday worship (\emph{First Apology} 65, 67). By the time represented in \emph{Ordo Romanus I} 12--15, the papal stational Mass includes an extensive reception of bread and wine from the people, a collection, mixing of the chalice, preparation of the altar, and offertory psalmody. That large urban action differed materially from a modern parish procession. Its evidence establishes deep Roman memories of presented gifts, diaconal service, charity, and song; it does not prove that every present route or object was restored from the seventh century.

Irenaeus supplies a patristic theology particularly apt for created gifts. Against views that despised material creation, he argues that the Church offers the Creator's own creatures in thanksgiving and that bread and cup, receiving God's invocation, bear an earthly and heavenly reality while promising the resurrection of embodied persons (\emph{Adversus haereses} IV.17--18; V.2). His polemical setting must be respected: he is not commenting on the present Offertory. Yet his reasoning explains why ordinary bread and wine matter. Redemption does not discard creation; Christ takes creation into the sacramental economy of his own Body and Blood.

The chant accompanies the procession and may continue at least until the gifts have been placed on the altar; the same norms governing the Entrance options apply as GIRM 74 specifies. It is not background sound for fundraising and not a second verbal commentary that must mention offering. Its ritual work is to give the corporate movement a common voice. When no procession occurs, chant may still accompany the preparation. Silence can expose the quiet formulas; singing can enlarge the common act. Neither condition changes the identity of the gifts or the purpose of the movement.

\begin{studybox}{What is offered at this point?}
Bread and wine are truly presented for the Eucharistic sacrifice, and the faithful join material alms and the offering of their lives. Yet the rite has not consecrated bread and wine, and Christian self-offering has no independent sacrificial efficacy. \emph{Sacrosanctum Concilium} 48 orders the relation: the faithful offer the immaculate Victim through the Priest's ministry and, with him, learn to offer themselves. Preparation is therefore neither empty staging nor a consecration before the Eucharistic Prayer.
\end{studybox}

\ordermovement{\latin{Benedictus es, Domine} over the Bread}{A Creature Received for Its Sacramental End}{Fixed quiet formula; it may be said aloud and answered when no Offertory chant is sung}{GIRM 75, 141; \latin{Ordo Missae}}

Standing at the altar, the Priest takes the paten with the bread, holds it slightly raised above the altar, and says the appointed formula quietly. He then places the paten upon the corporal. If no Offertory chant is being sung, the Missal permits him to say the formula aloud, after which the people may make the appointed acclamation (GIRM 141). The small lift is a presentation, not the elevation after the institution narrative, and the bread remains unconsecrated. Action, text, and sacramental state must be allowed to say different things.

The Latin incipit \latin{Benedictus es, Domine, Deus universi} places the utterance in the biblical world of blessing God for creation. Israel blesses the Lord who gives food and sustains life; 1~Chronicles~29 joins blessing, divine sovereignty, and the acknowledgment that every gift comes first from God; Psalm~104 praises the earth's fruitfulness and the bread that strengthens human life. The formula also names human cultivation. Grain does not arrive at the altar as untouched nature: earth, rain, seed, harvest, milling, baking, transport, and giving have converged in it. Human work is neither divine creation nor alien to it. It cooperates with created possibilities and is judged and healed within thanksgiving.

The exact prayer is modern. Its form is often described as a \emph{berakah}, and it genuinely adopts the scriptural-Jewish pattern of blessing the Creator over a gift. The postconciliar reform composed the Latin formula in a biblical blessing idiom to accompany the received bread-offering action. Ancient pattern, modern text, and current Roman use are three distinct claims.

Irenaeus again gives the most fitting patristic horizon. The Church's firstfruits do not supply something God lacks. They confess the goodness of creation, give thanks for divine bounty, and train the offerers against ingratitude (\emph{Adversus haereses} IV.17.5--18.6). Because the Creator and Redeemer are one, the bread's creaturely history can open toward its sacramental end. That end is not yet realized by the preparation formula. The words orient the gift toward the Eucharistic act in which Christ will give himself under this sign.

The formula's final horizon is therefore prospective. It does not merely thank God for ordinary nourishment, and it does not declare the bread already Christ's Body. It names what the bread is, whence it comes, how human hands have served it, and the liturgical destiny for which it has been received. The Priest's quiet voice ordinarily allows the common chant to remain the assembly's audible speech. When the formula is audible, the people's response ratifies the blessing. Both modes preserve the same ordered transition from received creation toward sacramental gift.

\ordermovement{The Mixed Chalice and \latin{Benedictus es, Domine} over the Wine}{Communion Written into the Cup}{Fixed mixture of water and wine with quiet prayer, then offering formula; the Deacon performs his assigned role when present}{GIRM 75, 142; \latin{Ordo Missae}}

At the side of the altar, the Deacon, or the Priest when there is no Deacon, pours wine and a little water into the chalice while saying the appointed quiet prayer. The chalice is then presented at the altar with a formula parallel to the blessing over bread, held slightly raised, and placed upon the corporal; it may be covered with a pall. As with the bread, the formula can be audible with the people's response when no chant is sung (GIRM 142). The rite thus contains two related but nonidentical units: the mixture with its private Christological petition, and the presentation of the prepared chalice with its blessing of the Creator.

A mixed cup is among the earliest securely attested Eucharistic actions. Justin twice speaks of water with wine in the cup brought to the president (\emph{First Apology} 65, 67). In the third century, Cyprian opposes the use of water alone and appeals to Christ's institution; he also reads the mingling typologically, taking the water as the people joined inseparably to Christ represented by the wine (\emph{Epistle} 63). His exegesis is a North African patristic interpretation rather than the only possible meaning of the act, but it reveals how early Christians saw ecclesial communion within material mixture.

The present private prayer gives the action an Incarnational and participatory horizon. Its language derives from an ancient Christmas collect on the wonderful creation and more wonderful restoration of human nature. That collect was adapted to the mingling of water and wine in the medieval West, probably by the eleventh century; the postconciliar Order retains an abbreviated form beginning \latin{Per huius aquae et vini mysterium}. The prayer's theological ancestry is therefore ancient, its assignment to this gesture medieval, and its present shortened form modern. The history exemplifies why the age of words, position, and ritual use must be dated separately.

The symbolism should not be pressed into material arithmetic. A ``little water'' does not quantify humanity's importance beside Christ, nor does the mixture itself effect the hypostatic union or incorporate the congregation sacramentally before consecration. It disposes the chalice according to immemorial Eucharistic use while the prayer asks that human participation in divine life be granted through the Incarnation. Cyprian's ecclesial reading and the current prayer's Christological reading are complementary horizons, not definitions of chemical parts.

The chalice formula beginning \latin{Benedictus es, Domine} is, like its bread counterpart, a postconciliar composition in biblical blessing form. Wine gathers creation and human labor differently from bread: vine, fruit, vintage, fermentation, craft, festivity, danger, and abundance all stand within the Bible's moral and eschatological range. Psalm~104 praises wine among God's gifts; Isaiah's banquet of rich wine images eschatological salvation (Isa.~25:6--9); Christ names the new fruit of the vine in the kingdom (Mark~14:25). The rite neither canonizes every human use of wine nor treats it as a neutral prop. It receives this created and worked gift for the cup of the new covenant.

\ordermovement{\latin{In spiritu humilitatis} and the Incensation}{Contrition and Fragrant Honor around One Offering}{Quiet priestly self-commendation fixed; gifts, cross, altar, Priest, and people may be incensed}{GIRM 75, 143--144}

After setting down the chalice, the Priest bows profoundly and says the quiet prayer beginning \latin{In spiritu humilitatis}. If incense is used, he then places it in the thurible and incenses the gifts, cross, and altar; the Deacon or another minister incenses the Priest and people (GIRM 143--144). The bow is fixed, the incensation optional. The two actions share the same threshold but should not be fused: one is the Priest's humble self-commendation before God; the other extends ordered signs of offering, prayer, and honor through the sanctuary and assembly.

The quiet prayer adapts the Vulgate form of Daniel~3:39--40, where Azariah, deprived of the Temple's sacrifices, approaches God with a contrite spirit and asks that humble persons be received. Its scriptural source is ancient; its use as a priestly Offertory prayer belongs to the medieval development of preparatory \emph{apologiae}. The postconciliar reform retained this compact biblical inheritance while removing other private offertory formulas. The prayer does not imply that the Eucharist lacks an objective sacrifice unless the Priest achieves a certain feeling. It places the minister personally beneath the mercy of the God before whom he is about to voice the Church's great thanksgiving.

Incense carries several biblical associations. Psalm~141:2 likens prayer to incense rising before God; Revelation~8:3--4 joins incense and the prayers of the saints at the heavenly altar; the Magi's gift and Temple worship add Christological and cultic resonances. These texts interpret fragrant smoke within Christian worship but do not date the present Offertory incensation. Early Roman evidence knows incense in other locations. \emph{Ordo Romanus I} includes it in the entrance and around the Gospel while giving no offertory incensation. Amalarius in the ninth century provides an early clear Western discussion of incensing the gifts while indicating that the usage was not then Roman; the practice entered and elaborated within the medieval Roman tradition.

The current rite simplifies that inheritance and supplies its own interpretation. Gifts, cross, and altar may be incensed so that the Church's offering and prayer rise before God; Priest and people are incensed because of their sacred dignity (GIRM 75). The route distinguishes without isolating. Bread and wine are honored as gifts destined for consecration, the cross names the once-for-all sacrifice, the altar is the sacramental table, the ordained minister will speak in his proper office, and the baptized people form the priestly body whose offering he voices. Incense does not make any of them consecrated matter, and honor does not make every recipient an equal liturgical agent.

Optionality matters. When incense is present, the sensory cloud can reveal an offering larger than the objects on the corporal: prayer and persons are drawn around Christ's altar. When absent, the rite is not less sacrificial and the people possess no lesser baptismal dignity. The optional sign interprets what the fixed order already contains; it does not supply a missing reality.

\ordermovement{\latin{Lavabo} (Washing of Hands)}{A Practical Gesture Turned toward Interior Cleansing}{Fixed priestly washing at the side of the altar with quiet prayer}{GIRM 76, 145}

The Priest goes to the side of the altar, where a minister pours water over his hands, and he quietly says the short prayer beginning \latin{Lava me, Domine} (GIRM 145). The action is ministerial and local: it is not a washing of every participant, a second Penitential Act, or a purification of the gifts. GIRM 76 states its present meaning as the Priest's desire for interior cleansing. That authoritative interpretation governs the movement even though its physical ancestry was more practical.

In the papal stational liturgy described by \emph{Ordo Romanus I} 14, the pontiff washes after receiving offerings from many hands. Cleanliness naturally belongs between that collection and the handling of bread, chalice, and altar. Ancient Eastern mystagogy could already disclose a moral meaning in a similar action. The Jerusalem catechist says that washing the hands signifies purification from sinful deeds before approaching the holy mysteries (\emph{Mystagogical Catechesis} 5.2). This is a valuable fourth-century interpretation, but it does not prove that the Roman handwashing descended from Jerusalem or that its original Roman motive was exclusively symbolic.

During the Middle Ages the Roman action acquired a longer recitation from Psalm~25/26 beginning \latin{Lavabo inter innocentes manus meas}. The present Order replaces that extended psalm with a brief petition drawn from Psalm~51:4. Thus the bodily action has an early practical Roman witness; explicit Christian moral interpretation is patristic and geographically broader; the former Roman text was medieval; and the current text and placement within the simplified preparation are products of reform. One movement holds all four histories without needing to pretend they began together.

The wash also clarifies ministerial symbolism. The Priest who presides remains a sinner dependent upon divine mercy. His hands will handle the gifts and the sacramental species because of ordination and liturgical office, not because they are morally spotless. Yet bodily cleanliness and inward conversion are not competitors. The practical body becomes capable of bearing prayer, and the quiet prayer keeps practical efficiency from exhausting the sign. At the altar's edge, water confesses that no minister approaches by self-certification.

\ordermovement{\latin{Orate, fratres} and the People's Response}{Distinct Ministries within the One Sacrifice}{Fixed public invitation and response; in the United States the people's standing begins here}{GIRM 77, 146}

Returning to the middle of the altar, the Priest faces the people and invites them with the formula beginning \latin{Orate, fratres}; the people stand and make the appointed response (GIRM 146 and the United States adaptation at GIRM 43). The Prayer over the Offerings follows. This is not a casual request for encouragement. It is the public hinge at which quiet preparation becomes explicit ecclesial consent to the sacrificial prayer about to begin.

The Latin phrase \latin{meum ac vestrum sacrificium} deliberately distinguishes ``mine'' and ``yours'' while the sentence directs both toward one acceptable sacrifice. The distinction is neither ownership nor rivalry. By Baptism, the faithful share in Christ's priestly people and offer spiritual sacrifices in him (1~Pet.~2:4--10). By ordination, the Priest acts in the person of Christ the Head and offers the Eucharistic sacrifice sacramentally in the Church's name. \emph{Lumen gentium} 10 teaches that common and ministerial priesthood differ essentially and not only in degree, while being ordered to one another. The invitation and response put that doctrine into pronouns.

\emph{Sacrosanctum Concilium} 48 completes the relation. The faithful are not silent outsiders watching the Priest's possession; through his hands they offer the immaculate Victim, and with him they learn to offer themselves. Nor does their participation make presidency a delegated function that the assembly could perform without ordination. Christ's sacrifice is one, the Church's participation is one, and the modes of participation remain differentiated. The exchange therefore resists both clerical isolation and an undifferentiated account of communal presidency.

Invitations to pray for the celebrant and sacrifice appear in the early medieval Western offertory complex, but the exact history of the current pair is gradual rather than apostolic. Forms of \latin{Orate, fratres} circulated with varying private responses; the now familiar invitation and popular answer are securely joined in the printed Roman Missal of 1474 and pass into the Missal of 1570. The postconciliar rite retains the inherited exchange, makes its public and communal character unmistakable, and locates the people's standing here in the United States. Biblical and patristic theology can explain its sacrificial communion; neither supplies its exact medieval words.

The response must be allowed to remain a response. It does not consecrate, substitute for the Eucharistic Prayer, or express a vote on the Priest's performance. It asks that the sacrifice be received for divine glory and the good of the Church, thereby widening personal intentions into an ecclesial horizon. The standing body then remains joined to the presidential prayer that seals preparation.

\ordermovement{\latin{Oratio super oblata} (Prayer over the Offerings)}{The Proper Petition That Seals Preparation}{One proper presidential prayer with a short conclusion and Amen}{GIRM 30, 32, 77, 146}

With hands extended, the Priest says the one appointed Prayer over the Offerings; the people answer Amen. Unlike the quiet preparatory formulas, this is a presidential prayer spoken aloud in the name of the holy people (GIRM 30, 32). No other prayer, song, or instrumental sound should compete with it. Its short conclusion and the people's ratification bring the preparation to completion and open directly into the Preface dialogue.

This prayer is proper: its exact theology belongs to the day, season, feast, ritual Mass, or intention whose formulary has been appointed. One text may ask that gifts be sanctified; another may interpret fasting, martyrdom, Paschal joy, nuptial covenant, or ecclesial service; another may ask that the offerers themselves be transformed. A generic exposition cannot silently substitute one composite meaning for all of them. What remains stable is the prayer's ritual location and relation: the gifts now on the altar are placed beneath the Church's petition before the Eucharistic Prayer begins.

Roman \emph{orationes super oblata} are among the ancient strata preserved in the Verona, Gelasian, and Gregorian sacramentary traditions. Their textual histories differ, and individual prayers must be dated on their own evidence. The later title \emph{Secreta} became associated with quiet recitation; a secure Frankish direction for the celebrant to pray quietly appears by the middle of the eighth century, while proposed etymologies and a single date of origin remain contested. In the Tridentine Order the prayer was said inaudibly and its conclusion supplied the transition to the Preface. The postconciliar reform restored the entire proper oration to audible presidential proclamation and the people's audible Amen.

That reform changed performance without inventing the prayer's office. Justin already attests that the president gives thanksgiving and the people ratify prayer with Amen, though he does not isolate a modern Prayer over the Offerings (\emph{First Apology} 65). Ancient Roman sacramentaries show a proper prayer over the presented gifts. The current arrangement reunites that Roman presidential text with a perceptible communal response. Antiquity of function, antiquity of many formularies, later quiet performance, and modern audibility must again be distinguished.

The Latin title \latin{super oblata} names gifts brought forward, not consecrated species. The prayer may use sacrificial language because bread and wine are destined for the Eucharistic sacrifice and because the whole Church is entering Christ's offering. It does not imply that consecration has already occurred or that the next prayer merely adds devotion to a completed oblation. The Amen seals preparation; it does not seal the Eucharistic memorial. Immediately the Preface dialogue lifts the Church into the great act of thanksgiving in which the gifts' sacramental identity will be effected and Christ's one sacrifice made present.
